% !TeX root = ../navier-stokes-manuscript.tex
\subsection{Where the two directions meet}\label{sub:ThePressureCompleteMixedDerivativeTower}

The vertical direction came from the singularity.
If velocity loses all \(L^3\) control before one finite time, then every fixed temporal response has to lose it too.
The horizontal direction comes from the equation that governs those responses.
Viscosity keeps two spatial derivatives of \(V_n\) in every temporal row, while transport and pressure keep the lower rows coupled to it.

Critical height makes the turn visible.
For
\[
  E_s(t):=\frac12\norm{\Lambda^su(t)}_2^2,
\]
the first differentiated balance has the form
\[
  H'
  =
  \mathcal P_{1/2}
  -2\nu E_{3/2},
\]
where \(\mathcal P_{1/2}\) retains the nonlinear and pressure transfer.
Differentiate again and the new highest viscous term is \(E_{5/2}\).
At order \(n\), it is \(E_{n+1/2}\), with all earlier transfer terms still attached.
Section~3 completes that calculation and isolates the top spatial height without pretending the other terms have vanished.

\Needspace{17\baselineskip}
Put only the positions of the two calculations beside one another:
\begin{center}
\captionsetup{hypcap=false}
\captionof{table}{The order shared by the temporal demand and the critical-height calculation.  The two middle entries are not the same object; their common order reveals the two directions inside the Tower.}
\label{tab:BodyTemporalSpatialDiscovery}
\small
\renewcommand{\arraystretch}{1.4}
\begin{tabularx}{\textwidth}{@{}c c >{\raggedright\arraybackslash}X c@{}}
\toprule
order & temporal response & critical-height response & highest spatial height exposed \\
\midrule
\(0\) & \(V_0=u\) & \(H=E_{1/2}\) & \(E_{1/2}\) \\
\(1\) & \(V_1=u_t\) & \(H'\) & \(E_{3/2}\) \\
\(2\) & \(V_2=u_{tt}\) & \(H''\) & \(E_{5/2}\) \\
\(3\) & \(V_3=u_{ttt}\) & \(H'''\) & \(E_{7/2}\) \\
\(n\) & \(V_n=\partial_t^nu\) & \(H^{(n)}\) & \(E_{n+1/2}\) \\
\bottomrule
\end{tabularx}
\end{center}

The table does not identify \(V_n\) with \(H^{(n)}\).
The scalar height derivatives are contractions of temporal rows, and they cannot reconstruct the velocity.
Nor does the appearance of \(E_{n+1/2}\) supply a bound.
It reveals the direction in which a bound would act.

That direction is the reversal.
The singularity asks for higher temporal responses to lose control.
The viscous calculation opens higher Sobolev levels in space.
One side seems to demand rougher behaviour, while the other exposes more spatial regularity, not less.
The two directions do not support the singularity in the same way.

\Needspace{12\baselineskip}
The mixed jet shows that this is not confined to the base row \(V_0=u\).
Fix any temporal order \(r\).
Its spatial coordinates run through
\[
  V_r\in H^0,\quad
  V_r\in H^1,\quad
  V_r\in H^2,\quad
  \ldots.
\]
At any time for which that row can be placed at every finite Sobolev level, the spaces meet as an intersection:
\[
  V_r(t)
  \in
  \bigcap_{m=0}^{\infty}H^m(\R^3)
  =
  H^\infty(\R^3)
  \subset
  C^\infty(\R^3).
\]
The spaces intersect.
Their norms are not summed, and no single constant is asked to control all orders at once.
The same intersection-to-smoothness implication is available separately in every temporal row.
This is the spatial meaning of the intersection, not yet a conclusion at \(T_*\).

\Needspace{15\baselineskip}
Lay those rows beneath one another and the full shape appears:
\[
\begin{array}{c|ccccc}
 & H^0_x & H^1_x & H^2_x & H^3_x & \cdots\\
\hline
V_0=u      & \bullet & \bullet & \bullet & \bullet & \cdots\\
V_1=u_t    & \bullet & \bullet & \bullet & \bullet & \cdots\\
V_2=u_{tt} & \bullet & \bullet & \bullet & \bullet & \cdots\\
V_3=u_{ttt}& \bullet & \bullet & \bullet & \bullet & \cdots\\
\vdots     & \vdots  & \vdots  & \vdots  & \vdots  & \ddots
\end{array}
\]
The downward coordinate is temporal order.
The across coordinate is spatial regularity.
Their meeting is already present in every entry of the pressure-complete mixed jet.

An estimate cannot take the infinite diagram as one indivisible object.
Choose a finite temporal depth \(N\) and a finite spatial depth \(M\).
The corresponding finite cut has the shape of a rectangle.
Section~3 builds the exact rectangle with the adjacent rows and pressure relations required by the equation, then proves the control that turns those finite cuts into an intersection.

The perfectly stretching vortex began with a motion that seemed able to press energy into smaller regions while its responses accelerated toward one finite time.
Following that motion opened the temporal Tower.
Following the equation inside the Tower opened the spatial direction.
Once the finite rectangles are controlled, the two directions meet in smoothness.
